\teachcover{Counting and building\\finite fields}{1}{
Bases turn algebra into counting
\[
\F_7\subset\F_{49}
\]
}{
This lecture finishes the vector-space reminder, proves the necessary
prime-power size theorem, and constructs the quadratic extension needed
for Fibonacci modulo $7$. The construction proves existence for this
example, not for every prime power. The main route now proves that
dimension is well defined by a six-slide exchange argument. Keep the
basis calculation break; the extra extension-arithmetic exercise is
optional material after the exit question.
}

\teachslide{Recall: scalars determine the combinations}{1}{
In $\F_3^2$, let $u=(1,2)$.
\[
\Span(u)=\{0u,1u,2u\}
=\{(0,0),(1,2),(2,1)\}.
\]
With $e_1=(1,0)$ and $e_2=(0,1)$,
\[
(a,b)=ae_1+be_2.
\]
\key{How can we tell whether coordinates are unique?}
}{
Retrieve rather than re-lecture: ask for $2u$, and ask how to express
$(2,1)$ in the standard coordinate vectors. A spanning list expresses
every vector, but could express the same vector in different ways.
The next two slides isolate that issue.
}

\teachslide{A redundant spanning list}{2}{
In $\F_3^2$, put $w=e_1+e_2$.

The list $e_1,e_2,w$ spans the space, but
\[
w=0e_1+0e_2+1w=1e_1+1e_2+0w.
\]
The difference gives a nontrivial relation:
\[
\key{e_1+e_2-w=0.}
\]
}{
Point out two distinct coefficient triples for the same vector.
Modulo $3$, the coefficients in the relation are $1,1,2$, so they
are not all zero. The problem is redundancy, not lack of spanning.
Ask which vector can be removed without changing the span.
}

\teachslide{Linear independence}{3}{
Vectors $v_1,\ldots,v_d$ are \key{linearly independent} if
\[
a_1v_1+\cdots+a_dv_d=0
\quad\Longrightarrow\quad a_1=\cdots=a_d=0.
\]
Then two expressions for the same vector have the same coefficients:
\[
\sum a_iv_i=\sum b_iv_i
\ \Longrightarrow\ \sum(a_i-b_i)v_i=0.
\]
}{
Finish the second line verbally: independence forces every $a_i-b_i=0$.
Conversely, unique expressions of zero force independence. Give the
standard pair $e_1,e_2$ as the basic example: $(a,b)=0$ implies both
coordinates are zero. Distinguish a nontrivial relation from one whose
coefficients all vanish.
}

\teachslide{A basis gives existence and uniqueness}{2}{
A \key{basis} is a linearly independent spanning list.

Every $v\in V$ has a unique expression
\[
\key{v=a_1v_1+\cdots+a_dv_d.}
\]
\begin{itemize}
\item Spanning gives existence of the coefficients.
\item Independence gives uniqueness.
\end{itemize}
}{
Have students supply which hypothesis does which job. An ordered basis
also specifies the order of coordinates. The standard basis of $k^d$
is the list of coordinate vectors. We will count ordered bases when
we count matrices in Lecture 3.
}

\teachslide{Build a basis one vector at a time}{3}{
Start with no vectors; their span is $\{0\}$.

If the current span is not all of $V$, choose a vector outside it.

Adding that vector preserves independence.
\spacer
\key{In a finite space, this process must stop.}

At stopping, the chosen vectors form a basis.
}{
Prove the key step. If $a_1v_1+\cdots+a_rv_r+bw=0$ and $b\ne0$,
then division by $b$ expresses $w$ in the old span, contradicting its
choice. Thus $b=0$, and old independence makes all other coefficients
zero. Finiteness makes the strictly increasing sequence of spans stop.
This proof specifically establishes the finite basis needed below.
}

\teachslide{Could two bases have different lengths?}{1}{
A basis gives unique coordinates.

But before defining dimension, we must prove that its length
does not depend on the chosen basis.
\spacer
Our key comparison:
\[
\key{\text{independent list length}\ \le\ \text{spanning list length}.}
\]
}{
Make the logical issue explicit. Existence and uniqueness of
coordinates refer to one chosen basis. They do not yet compare two
different bases. The following argument works over any field,
including infinite fields; it does not use the count $p^d$.
Only finite lists are involved in the comparison theorem.
}

\teachslide{Exchange one vector}{2}{
Suppose $w_1,\ldots,w_m$ span $V$ and $v\ne0$.
\[
v=a_1w_1+\cdots+a_mw_m.
\]
Some coefficient is nonzero; relabel so $a_1\ne0$.
\[
\key{w_1=a_1^{-1}\bigl(v-a_2w_2-\cdots-a_mw_m\bigr).}
\]
Thus $v,w_2,\ldots,w_m$ still span $V$.
}{
Explain both inclusions of spans. Each vector in the new list lies
in the old span. The solved equation puts the discarded $w_1$ in
the new span, and all other old vectors remain. Hence the spans are
equal. The spanning list need not be independent. Relabeling just
moves a nonzero coefficient to the first position. Dividing by it
is valid because the scalars form a field.
}

\teachslide{An exchange in $\F_3^2$}{2}{
Start with $e_1=(1,0)$ and $e_2=(0,1)$.

Insert $u=(1,1)=e_1+e_2$.
\[
e_1=u-e_2,\qquad \key{u,e_2\text{ still span}.}
\]
Now insert $v=(1,2)=u+e_2$.
\[
e_2=v-u,\qquad \key{u,v\text{ still span}.}
\]
}{
Ask students to supply each subtraction. Over $\F_3$ a minus sign
means adding twice that vector. Each exchange preserves the number
of vectors and the span. The general proof will also need to ensure
we can always remove an old $w$, rather than a vector already inserted
from the independent list. This example anticipates the retained
basis exercise later in the lecture.
}

\teachslide{Repeat the exchange}{3}{
Let $v_1,\ldots,v_r$ be independent and $w_1,\ldots,w_m$ span $V$.

After inserting $v_1,\ldots,v_k$, the list still spans and has
\[
v_1,\ldots,v_k,\quad m-k\text{ remaining old vectors}.
\]
Express $v_{k+1}$ in this list.

\key{Some remaining old vector has a nonzero coefficient.}

Otherwise $v_{k+1}\in\Span(v_1,\ldots,v_k)$, contradicting independence.
}{
Start with $k=0$. Since the list spans, such an expression exists.
If all coefficients of the remaining old vectors vanished, moving
$v_{k+1}$ to the other side would give a relation among the $v_i$
with its coefficient -1, contradicting independence. Therefore
solve for an old vector having nonzero coefficient and replace that
vector by $v_{k+1}$. The earlier exchange formula preserves spanning
and never removes an already inserted $v_i$.
}

\teachslide{An independent list cannot be longer}{2}{
Every insertion replaces one remaining old vector.

After $m$ insertions, the spanning list would be
\[
v_1,\ldots,v_m.
\]
If $r>m$, then $v_{m+1}$ would be in their span, contradicting independence.
\[
\key{r\le m.}
\]
This is the \key{Steinitz exchange argument}.
}{
The contradiction is the rigorous reason the process cannot outlast
the original spanning list. If fewer than $m$ independent vectors
were given, the process stops earlier and the inequality already
holds. If $m=0$, the space is $\{0\}$ and any nonempty independent
list is impossible, so the assertion remains true. The argument
also shows an independent finite list can be extended to a basis
when a finite spanning list is available: keep the inserted vectors
and remove any remaining redundancies.
}

\teachslide{Dimension is well defined}{2}{
Suppose $B$ and $C$ are bases, of lengths $r$ and $m$.
\begin{itemize}
\item $B$ independent, $C$ spanning: $r\le m$.
\item $C$ independent, $B$ spanning: $m\le r$.
\end{itemize}
Therefore \key{$r=m$}.

For a space with a finite basis, this common length is its dimension:
\[
\dim_k V=r.
\]
}{
Only now define dimension. Such a space is called finite-dimensional.
The result proves independence from the choice of basis, not
independence from the choice of scalar field. The theorem compares
finite bases. It also rules out an infinite independent basis when
a finite spanning list exists, since an arbitrarily long finite
sublist would violate the inequality. No Gaussian elimination was
needed, though elimination supplies another proof of the same
comparison theorem. The zero space has the empty basis and dimension
zero. The next slide specializes to finite-field coordinates.
}


\teachslide{Count coordinate choices}{3}{
Suppose $v_1,\ldots,v_d$ is a basis over $\F_p$.
\[
v=a_1v_1+\cdots+a_dv_d,
\qquad a_i\in\F_p.
\]
Each of the $d$ coordinates has $p$ choices.

Uniqueness means that different choices give different vectors.
\[
\key{|V|=p^d.}
\]
}{
The coordinate map from $\F_p^d$ to $V$ is onto by spanning and
one-to-one by independence. Thus the product rule counts the vectors
exactly, not merely bounds their number. In particular, a span of
$r$ independent vectors has $p^r$ elements. If two bases have lengths
$d,e$, this count gives $p^d=p^e$ and hence $d=e$.
}

\teachslide{Your turn: is this a basis?}{3}{
In $\F_3^2$, let $u=(1,1)$ and $v=(1,2)$.
\begin{enumerate}
\item Solve $au+bv=(0,0)$.
\item Express $(2,0)$ as $au+bv$.
\item How many vectors are in $\Span(u,v)$?
\end{enumerate}
\spacer
\key{Write the two coordinate equations.}
}{
Give time to write $a+b=0$ and $a+2b=0$. Subtract the first from
the second to get $b=0$, then $a=0$. For $(2,0)$, use $a+b=2$,
$a+2b=0$. A suggestion for students stuck in real-number habits:
$-2=1$ in $\F_3$. The answer appears on the next page.
}

\teachslide{Basis exercise: answers}{1}{
\[
a+b=0,\quad a+2b=0\quad\Longrightarrow\quad b=0, a=0.
\]
So $u,v$ are independent.
\[
\key{(2,0)=u+v.}
\]
Their span has $3^2=9$ elements, all of $\F_3^2$.

Thus $u,v$ form a basis.
}{
There are nine elements in the ambient space, so the nine-element
span is the whole space. We have checked both requirements for a
basis without silently assuming that independence implies spanning
in arbitrary spaces. Coordinates depend on the chosen basis:
$(2,0)$ has coordinates $(1,1)$ in this basis.
}

\teachslide{A field can be a vector space}{3}{
Let $K$ be a finite field of characteristic $p$.

It contains $\F_p$. Regard
\begin{itemize}
\item elements of $K$ as vectors;
\item elements of $\F_p$ as scalars;
\item multiplication in $K$ as scalar multiplication.
\end{itemize}
\[
a(u+v)=au+av,\qquad a(bu)=(ab)u.
\]
}{
Check all the vector-space axioms, not just the two displayed ones:
addition in $K$ is an abelian group operation, $(a+b)u=au+bu$, and
$1u=u$. These are field laws restricted to the scalar subfield.
This is not making $K$ a vector space over the real numbers.
Here a vector can also be multiplied by another vector because $K$
has additional field structure; the vector-space definition alone
does not supply that multiplication.
}

\teachslide{Every finite field has prime-power size}{2}{
Finiteness gives a basis of $K$ over its prime field $\F_p$.
\[
\key{|K|=p^d,\qquad p\text{ prime},\quad d\ge1.}
\]
\begin{itemize}
\item Characteristic gives the prime $p$.
\item A basis over $\F_p$ gives the exponent $d$.
\end{itemize}
\spacer
There is no field with $6$ elements.

A field with $49$ elements would have characteristic $7$.
}{
Ask students to rule out $6$ as a prime power. For $49$, unique prime
factorization gives $p=7,d=2$. We have proved a necessary condition;
we have not proved existence or uniqueness for arbitrary $p^d$.
The next construction supplies a particular missing example.
To justify a finite basis, apply the earlier basis-building process
to the finite set $K$. The strictly increasing spans must stop.
Since $1\ne0$, its dimension over $\F_p$ is at least one. Thus
finiteness first gave prime characteristic and now gives a finite basis.
}

\teachslide{Build the missing square root}{2}{
Over $\F_7$, $5$ is not a square.

Introduce a symbol $s$ with $s^2=5$, and use formal pairs
\[
\key{K=\{a+bs:a,b\in\F_7\}.}
\]
Two elements are equal exactly when their two coefficients agree.
\[
(a+bs)+(c+es)=(a+c)+(b+e)s.
\]
}{
Define the elements as pairs first, so uniqueness of this representation
is built in. The symbol $s$ is not an ordinary real square root.
We are defining new finite arithmetic. Write $s=(0,1)$ and embed
$a\in\F_7$ as $(a,0)$. There are already $7^2$ candidate elements;
we still must prove this arithmetic is a field.
}

\teachslide{Multiply, then replace $s^2$ by $5$}{3}{
\[
(a+bs)(c+es)=(ac+5be)+(ae+bc)s.
\]
For example,
\[
\begin{aligned}
(2+s)(3+2s)&=6+7s+2s^2\\
&=16+7s\\
&=2.
\end{aligned}
\]
All coefficients are reduced modulo $7$.
}{
Explain the four terms from distribution. This is polynomial arithmetic
with the relation $s^2=5$; the ring laws hold. If requested, associativity
can be checked from the pair formula, or justified by reducing products
of polynomials modulo $X^2-5$. A backup slide introduces that notation.
The product $2$ is an element of the embedded $\F_7$.
}

\teachslide{Why every nonzero element has an inverse}{3}{
\[
(a+bs)(a-bs)=a^2-5b^2.
\]
For $(a,b)\ne(0,0)$, this denominator is nonzero:
\begin{itemize}
\item If $b=0$, then $a^2\ne0$.
\item If $b\ne0$, its vanishing would give $(a/b)^2=5$.
\end{itemize}
Therefore
\[
\key{(a+bs)^{-1}=\frac{a-bs}{a^2-5b^2}.}
\]
}{
The second case contradicts that $5$ is nonsquare in $\F_7$. The
denominator is a nonzero element of the original field, so its inverse
exists there. This proves the new ring is a field. Explain the division
notation as multiplication by that inverse. If $5$ were a square, the
same formal construction could have nonzero zero divisors; a backup
slide gives the example modulo $11$.
}

\teachslide{A field with forty-nine elements}{2}{
We have built a field $K$, which we call $\F_{49}$.
\[
\key{\{1,s\}\text{ is a basis over }\F_7.}
\]
Every element is a unique $a\cdot1+b\cdot s$.
\[
|K|=7^2=49,\qquad \dim_{\F_7}K=2.
\]
Its characteristic is $7$.
}{
Spanning follows from the definition as pairs, and independence follows
because $a+bs=0$ means both coordinates are zero. Seven copies of $1$
sum to zero; no smaller positive number does. Thus the characteristic
is $7$, although the number of elements is $49$. The same construction
works over an odd prime with any nonsquare $d$ replacing $5$.
}

\teachslide{$\F_{49}$ and $\Z/49\Z$ are different}{2}{
\[
\begin{array}{c|c|c}
&\F_{49}&\Z/49\Z\\\hline
\text{Elements}&49&49\\
\text{Characteristic}&7&49\\
\text{Every nonzero element a unit?}&\text{Yes}&\text{No}
\end{array}
\]
In $\Z/49\Z$, the nonzero residue $7$ satisfies $7^2=0$.

In $\F_{49}$, seven copies of $1$ already give $0$.
}{
Read the two columns slowly. The symbol $7$ denotes different things
in the two arithmetics: a nonzero residue in one and the sum $7\cdot1=0$
in the other. Equal cardinality is far weaker than equal arithmetic.
This is the main misconception to address before moving on.
}

\teachslide{Exit question}{2}{
Suppose $K$ is a field with $25$ elements.
\begin{enumerate}
\item What is its characteristic?
\item What is its dimension over its prime field?
\item Why is $K$ not $\Z/25\Z$?
\end{enumerate}
}{
Expected answers: characteristic $5$, dimension $2$ over $\F_5$, and
$\Z/25\Z$ has characteristic $25$ and the nonzero zero divisor $5$.
The question assumes a field with $25$ elements; it does not rely on
a general finite-field existence theorem. This is the 50-minute stop.
}

\backups
\backupslide{Polynomial notation for the construction}{
\[
\key{\F_{49}=\F_7[X]/(X^2-5).}
\]
Two polynomials represent the same element when their difference
is a multiple of $X^2-5$.

Polynomial division gives a unique remainder $a+bX$.

The class of $X$ is our element $s$, and $s^2=5$.
}{
This is the polynomial analogue of integer residues. Addition and
multiplication descend from polynomial arithmetic, so the ring laws
follow. Polynomial division by the monic quadratic gives unique
remainders of degree at most one. The main lecture proves the field
property directly, so a general quotient-field theorem is unnecessary.
}

\backupslide{What if the chosen number is already a square?}{
Try the same formal-pair construction over $\F_{11}$, with $s^2=5$.

Since $4^2=5$,
\[
(s-4)(s+4)=s^2-16=0.
\]
Both factors are nonzero formal pairs.
\spacer
\key{This construction gives a ring with zero divisors, not a field.}
}{
The coefficient of $s$ is $1$ in both factors, so neither is the
zero pair. This makes the nonsquare hypothesis in the inverse proof
concrete. There is no need to construct an extension when the roots
already lie in the original field.
}

\backupslide{Your turn: calculate in the new arithmetic}{
In $K$, where $s^2=5$ and coefficients lie in $\F_7$, compute
\[
(1+2s)(3+s).
\]
Then compute
\[
(2+s)(2-s).
\]
\key{Distribute first; reduce afterward.}
}{
Allow a minute of quiet work. Ask students to show intermediate
expressions so that an error in reduction can be distinguished from
an error in expansion. The second calculation prepares the conjugate
inverse formula and should not be described as complex conjugation.
}


\backupslide{Two products}{
\[
(1+2s)(3+s)=3+7s+2s^2=13=6.
\]
\[
\key{(2+s)(2-s)=4-s^2=4-5=6.}
\]
Since $6^{-1}=6$ in $\F_7$,
\[
(2+s)^{-1}=6(2-s)=5+s.
\]
}{
Verify $6\cdot6=1\pmod7$. Multiplying the second identity by $6$
gives the inverse of this one element. Do not yet claim that every
element is invertible; the next slide proves that general statement.
For a quick check, $(2+s)(5+s)=15+7s=1$.
}


\backupslide{The roots now exist}{
In $\F_{49}$, division by $2$ is allowed, and
\[
\alpha=\frac{1+s}{2},\qquad\beta=\frac{1-s}{2}.
\]
\[
\alpha+\beta=1,\qquad\alpha\beta=-1.
\]
Both solve $X^2-X-1=0$.
\spacer
Next: \key{multiplicative orders, matrices, and the period.}
}{
Modulo $7$, $2^{-1}=4$, so $\alpha=4+4s$ and $\beta=4-4s$.
Their sum and product verify the factorization
$(X-\alpha)(X-\beta)=X^2-X-1$. The roots are distinct because $s\ne0$.
This reconnects the construction to the original motivating question.
Lecture 3 supplies cyclicity and matrix counts; Lecture 4 uses the
roots' powers to finish the period argument.
}


\backupslide{A bridge to matrices}{
An invertible matrix has columns that form an ordered basis.
\[
\key{\text{Counting invertible matrices = counting ordered bases.}}
\]
We will count the column choices one at a time.
}{
Ask students to remember how many elements are in the span of one
nonzero vector over $\F_p$: exactly $p$. That is the key number for
the next lecture's second column. This optional bridge previews
Lecture 3; the main route already ended at the exit question.
}

\backupslide{Gaussian elimination gives another proof}{
Suppose $w_1,\ldots,w_m$ span and $v_1,\ldots,v_r$ are independent.

Choose coefficients $v_i=\sum_{j=1}^m c_{ji}w_j$; let $C=(c_{ji})$.

If $r>m$, elimination in the homogeneous system $Cx=0$
leaves a free variable, hence a nonzero solution.
\[
\sum_{i=1}^r x_iv_i=\sum_{j=1}^m(Cx)_jw_j=0,
\]
contradicting independence.
}{
$C$ has $m$ rows and $r$ columns. Elimination produces at most $m$
pivots, so when $r>m$ a variable is free. Set one free variable to 1,
the others to 0, and solve backwards for the pivot variables. Row
operations preserve the solutions, and division by a nonzero pivot
is allowed over a field. The coefficients $c_{ji}$ need not be unique:
the original $w_j$ merely span, and any chosen expressions suffice.
Thus this proof does not assume the dimension theorem or rank-nullity.
It is an alternative to the main exchange proof, not a prerequisite.
}

